\documentclass[11pt,a4paper]{article}
\usepackage[T1]{fontenc}
\usepackage[utf8]{inputenc}
\usepackage{amsmath,amssymb,graphicx,hyperref,longtable,booktabs,array}
\usepackage[margin=2.2cm]{geometry}
\providecommand{\tightlist}{\setlength{\itemsep}{0pt}\setlength{\parskip}{0pt}}
\setlength{\parindent}{0pt}\setlength{\parskip}{6pt}

\begin{document}
\hypertarget{physics-informed-lag-compensation-for-continuous-glucose-monitoring}{%
\section{Physics-informed lag compensation for continuous glucose
monitoring}\label{physics-informed-lag-compensation-for-continuous-glucose-monitoring}}

DH 803 (Wearable Health Technology) group project, group \textbf{Lag
Compensators}: Harsh Jadhav (Chemical Engineering) and Shreyas Kamble
(Mechanical Engineering).

CGM sensors measure glucose in interstitial fluid, which lags blood
glucose by roughly 5 to 15 minutes. This project models that lag as
first-order transport and asks two questions in simulation:

\begin{enumerate}
\def\labelenumi{\arabic{enumi}.}
\tightlist
\item
  How much of the lag error can a physics-based state estimator (Kalman
  filter, RTS smoother) remove, in real time and retrospectively,
  compared with simple baselines?
\item
  Can the lag time constant tau be identified from data, per patient or
  for a population, and can physiological lag be separated from sensor
  lag?
\end{enumerate}

Status: simulation study complete through Piece 5 (31 Oct 2026). No
real-patient data yet; OhioT1DM validation needs a data use agreement.

\begin{quote}
\textbf{Attribution.} The code, analyses and notes in this repository
were written by Claude (an AI assistant, Anthropic) at the request of
Harsh Jadhav, across two chat sessions. The team has not yet
independently re-run or reviewed them. Before any part is presented as
the team\textquotesingle s own work, read it, re-run it, and record who
checked what. Commits are authored as "Claude" for this reason.
\end{quote}

\hypertarget{main-results-simulation-27-virtual-patients-5-fold-subject-wise-cross-validation}{%
\subsection{Main results (simulation, 27 virtual patients, 5-fold
subject-wise
cross-validation)}\label{main-results-simulation-27-virtual-patients-5-fold-subject-wise-cross-validation}}

Error during rapid glucose change (\textbar rate\textbar{} \textgreater=
1 mg/dL/min), scenario S0 (model matches), RMSE in mg/dL, median over
patients:

\begin{longtable}[]{@{}llll@{}}
\toprule\noalign{}
Method & Uses future data? & Noise sd 2 & Noise sd 8 \\
\midrule\noalign{}
\endhead
\bottomrule\noalign{}
\endlastfoot
Raw CGM & no & 19.7 & 21.5 \\
Naive inversion, tau = 10 & no & 10.2 & 17.9 \\
Kalman filter, tau = 10 & no & 7.5 & 18.4 \\
Time shift (tuned) & yes & 8.0 & 11.3 \\
RTS smoother, tau = 10 & yes & 4.4 & 8.4 \\
\end{longtable}

Key findings:

\begin{itemize}
\tightlist
\item
  \textbf{Real-time compensation is noise-limited.} The Kalman filter
  matches the future-peeking shift at low noise but does not beat naive
  inversion at noise sd 8, and earlier hypoglycaemia alerts come with
  more false alarms.
\item
  \textbf{Retrospective smoothing works best.} The smoother beats the
  tuned time shift (S0, noise 8: -1.9 mg/dL, 95\% CI -3.3 to -1.3) and
  cuts false alarms from 0.93 to 0.23 per day.
\item
  \textbf{Knowing the true tau barely helps the real-time filter}
  (difference about 0.1 mg/dL); noise, not tau uncertainty, is the
  bottleneck.
\item
  \textbf{Per-patient tau from the model likelihood is confidently
  wrong.} 95\% intervals cover the truth in only 4 to 19\% of patients,
  because the generic glucose model is misspecified for meal-driven
  glucose.
\item
  \textbf{A population lag chosen from fingersticks works.} Using about
  264 training fingersticks (no ground truth), it picked tau = 16 min
  when the sensor added 5 min of unmodelled lag, cutting smoother error
  from 13.5 to 5.8 mg/dL (true-lag oracle 4.5).
\item
  \textbf{Physiological and sensor lag cannot be separated} from these
  data; only their sum (an effective lag) is identifiable.
\end{itemize}

Figures: \href{results/fig_scenarios.png}{scenarios},
\href{results/fig_baselines.png}{baselines},
\href{results/fig_kalman.png}{Kalman estimators},
\href{results/fig_ident.png}{identifiability}.

\hypertarget{repository-layout}{%
\subsection{Repository layout}\label{repository-layout}}

\begin{verbatim}
code/        sensor layer, evaluation harness, Kalman/smoother, identifiability (modules, tests, run and plot scripts)
data/        sim.npz: 30 simglucose virtual patients x 5 days, 1-min plasma and subcutaneous glucose
results/     result JSON files, figures, split manifest, likelihood profiles
prototype/   first prototype (EKF with online tau) from chat 1, with its results and figures
docs/        project brief, prototype handoff, one note per piece (methods, tests, numbers, caveats), Assignment 2 deck
run_all.sh   runs every test and experiment in order
\end{verbatim}

\begin{longtable}[]{@{}llll@{}}
\toprule\noalign{}
Piece & What it adds & Code & Note \\
\midrule\noalign{}
\endhead
\bottomrule\noalign{}
\endlastfoot
1 & Verification of the prototype & \texttt{prototype/} &
\href{docs/piece-1-prototype-verification.md}{piece-1} \\
2 & Realistic sensor layer, scenarios S0 to S5 & \texttt{sensor.py} &
\href{docs/piece-2-sensor-layer.md}{piece-2} \\
3 & Evaluation harness, baselines B0 to B3 & \texttt{harness.py} &
\href{docs/piece-3-evaluation-harness.md}{piece-3} \\
4 & Kalman filter, RTS smoother, regularised inversion &
\texttt{kalman.py} &
\href{docs/piece-4-kalman-estimators.md}{piece-4} \\
5 & Identifiability of tau, fingerstick tuning & \texttt{ident.py} &
\href{docs/piece-5-identifiability.md}{piece-5} \\
\end{longtable}

Scenarios: S0 matched first-order lag; S1 extra 3-min sensor filter and
2-min delay; S2 two-stage lag; S3 tau drifting about 25\%; S4
calibration drift and compression lows; S5 all of these together.

\hypertarget{how-to-run}{%
\subsection{How to run}\label{how-to-run}}

Python 3.10 or newer.

\begin{verbatim}
pip install -r requirements.txt
bash run_all.sh            # tests (about 1 min) then all experiments (about 15 min on a laptop)
\end{verbatim}

Or step by step from \texttt{code/}:

\begin{verbatim}
cd code
python3 test_sensor.py ../data/sim.npz
python3 test_harness.py ../data/sim.npz
python3 test_kalman.py ../data/sim.npz
python3 test_ident.py
python3 run_kalman.py ../data/sim.npz ../results && python3 plot_kalman.py ../results/kalman_results.json ../results/fig_kalman.png
\end{verbatim}

\texttt{data/sim.npz} can be regenerated with
\texttt{cd\ prototype\ \&\&\ python3\ run\_gen.py} (needs
\texttt{simglucose}, about 7 min). That writes to \texttt{data/sim.npz}
and overwrites the shipped copy.

\hypertarget{caveats}{%
\subsection{Caveats}\label{caveats}}

\begin{itemize}
\tightlist
\item
  Simulation only. In simglucose the lag is exactly first-order, so
  matched-scenario results are best-case; scenarios S1 to S5 add
  deliberate mismatch.
\item
  Tuned time-shift and "population tau from truth" results use
  ground-truth glucose and are simulation-only references.
\item
  Sensor noise levels are anchored loosely to one published accuracy
  study; noise sd 8 mg/dL is the headline stress level.
\item
  Hypoglycaemia is over-represented in the simulated patients; three
  virtual children were excluded as unphysiological.
\end{itemize}

\hypertarget{next-steps}{%
\subsection{Next steps}\label{next-steps}}

\begin{itemize}
\tightlist
\item
  Piece 6 (optional): meal- and insulin-driven glucose model as the
  prior, to fix the likelihood misspecification.
\item
  Piece 7: OhioT1DM data use agreement and exploratory real-data check
  with fingersticks as reference.
\item
  Piece 8: ablations (AR(1) noise state, posterior coverage), final
  report.
\end{itemize}

\hypertarget{key-references}{%
\subsection{Key references}\label{key-references}}

\begin{enumerate}
\def\labelenumi{\arabic{enumi}.}
\tightlist
\item
  A. Basu et al., "Time lag of glucose from intravascular to
  interstitial compartment in humans," \emph{Diabetes}, vol. 62, no. 12,
  pp. 4083-4087, 2013. \url{https://doi.org/10.2337/db13-1132}
\item
  G. Schmelzeisen-Redeker et al., "Time delay of CGM sensors: relevance,
  causes, and countermeasures," \emph{J. Diabetes Sci. Technol.}, vol.
  9, no. 5, pp. 1006-1015, 2015.
  \url{https://doi.org/10.1177/1932296815590154}
\item
  K. Rebrin, G. M. Steil, W. P. van Antwerp, J. J. Mastrototaro,
  "Subcutaneous glucose predicts plasma glucose independent of insulin,"
  \emph{Am. J. Physiol. Endocrinol. Metab.}, vol. 277, no. 3, pp.
  E561-E571, 1999. \url{https://doi.org/10.1152/ajpendo.1999.277.3.e561}
\item
  B. W. Bequette, "Continuous glucose monitoring: real-time algorithms
  for calibration, filtering, and alarms," \emph{J. Diabetes Sci.
  Technol.}, vol. 4, no. 2, pp. 404-418, 2010.
  \url{https://doi.org/10.1177/193229681000400222}
\item
  A. Facchinetti et al., "Modeling the glucose sensor error," \emph{IEEE
  Trans. Biomed. Eng.}, vol. 61, no. 3, pp. 620-629, 2014.
  \url{https://doi.org/10.1109/TBME.2013.2284023}
\item
  C. Dalla Man et al., "The UVA/PADOVA type 1 diabetes simulator: new
  features," \emph{J. Diabetes Sci. Technol.}, vol. 8, no. 1, pp. 26-34,
  2014. \url{https://doi.org/10.1177/1932296813514502}
\item
  J. Xie, "simglucose v0.2.1," 2018.
  \url{https://github.com/jxx123/simglucose}
\item
  C. Marling and R. Bunescu, "The OhioT1DM dataset for blood glucose
  level prediction: update 2020," \emph{CEUR Workshop Proc.}, vol. 2675,
  2020. \url{https://pmc.ncbi.nlm.nih.gov/articles/PMC7881904/}
\end{enumerate}

This is a research and education project. Nothing here is a medical
device or a basis for treatment decisions.
\end{document}
